\documentclass[../cours_2de.tex]{subfiles}
\begin{document}
\setcounter{chapter}{6}
\chapter{Repérage, distances et lieux géométriques}
\label{sec:ch07}
\begin{objectif}
Passer d'une figure à des coordonnées et démontrer des propriétés par le calcul. Prérequis : Pythagore, milieu, médiatrice, repérage dans le plan.
\end{objectif}
\section{Choisir un repère}
Un repère du plan comporte une origine O et deux directions graduées non parallèles. Les coordonnées $(x;y)$ indiquent les déplacements selon ces directions. Un repère \textbf{orthogonal} a des axes perpendiculaires. Il est \textbf{orthonormé} si, en plus, les unités des deux axes ont la même longueur.

Les calculs de coordonnées des milieux se font dans tout repère. La formule usuelle de distance ci-dessous exige un repère orthonormé. Choisir l'origine sur un sommet et les axes sur des côtés perpendiculaires peut simplifier un problème.
\section{Distance entre deux points}
\begin{theoreme}
Dans un repère orthonormé, pour $A(x_A;y_A)$ et $B(x_B;y_B)$,
\[
AB=\sqrt{(x_B-x_A)^2+(y_B-y_A)^2}.
\]
\end{theoreme}
\begin{demonstration}
Posons $H(x_B;y_A)$. Les directions des axes donnent un angle droit en H. Les longueurs des deux côtés sont $AH=|x_B-x_A|$ et $HB=|y_B-y_A|$. Pythagore donne $AB^2=AH^2+HB^2$, puis la racine positive donne AB. Si les points ont la même abscisse ou la même ordonnée, l'un des côtés est nul et la formule se réduit à la distance sur un axe ; elle reste valide.
\end{demonstration}
\begin{center}
\begin{tikzpicture}[scale=.8]
\draw[->] (-.3,0)--(6,0) node[right] {$x$};
\draw[->] (0,-.3)--(0,4.5) node[above] {$y$};
\coordinate (A) at (1,1);\coordinate (H) at (5,1);\coordinate (B) at (5,4);
\draw[thick,accent] (A)--(B);\draw[dashed] (A)--(H)--(B);
\draw (4.75,1)--(4.75,1.25)--(5,1.25);
\fill (A) circle(2pt) node[left] {$A$};
\fill (B) circle(2pt) node[right] {$B$};
\fill (H) circle(2pt) node[right] {$H$};
\node[below] at (3,1) {$|x_B-x_A|$};
\end{tikzpicture}
\end{center}
Le théorème de Pythagore et sa réciproque sont démontrés dans le cours de troisième, chapitre 13 ; on les utilise ici comme acquis, sans prétendre les déduire de la formule de distance qui en dépend.
\begin{exemple}
Pour $A(-1;2)$ et $B(5;10)$, $AB=\sqrt{6^2+8^2}=10$. Comparer les carrés des distances évite souvent des racines inutiles : un triangle de côtés de carrés $9,16,25$ est rectangle par la réciproque de Pythagore.
\end{exemple}
\section{Coordonnées d'un milieu}
\begin{propriete}
Le milieu I du segment AB a pour coordonnées
\[
I\left(\frac{x_A+x_B}{2};\frac{y_A+y_B}{2}\right).
\]
\end{propriete}
\begin{demonstration}
Le déplacement de A vers I est la moitié de celui de A vers B. Sur le premier axe, $x_I-x_A=(x_B-x_A)/2$, d'où la formule pour $x_I$ ; même raisonnement pour y. Réciproquement, ces deux égalités donnent le même demi-déplacement sur chaque axe : I est sur AB et partage le segment en deux parties égales. La formalisation vectorielle est donnée au chapitre 8.
\end{demonstration}
Pour les points précédents, $I(2;6)$. Les différences divisées par deux et les sommes divisées par deux ne jouent pas le même rôle : le milieu est un point, pas seulement un déplacement.
\section{Cercles, disques et médiatrices}
Un cercle de centre $C(a;b)$ et de rayon $r>0$ est l'ensemble des points à distance r de C. Le disque fermé contient ceux dont la distance à C est au plus r.
\begin{propriete}
Le cercle a pour équation $(x-a)^2+(y-b)^2=r^2$ ; le disque fermé correspond à l'inégalité $\le r^2$.
\end{propriete}
\begin{demonstration}
La formule de distance donne $CM=\sqrt{(x-a)^2+(y-b)^2}$. Les distances et r sont non négatifs ; les comparer équivaut à comparer leurs carrés.
\end{demonstration}
La médiatrice de AB, pour $A\ne B$, est l'ensemble des points équidistants de A et B. Dans un repère orthonormé d'origine au milieu I et d'axe des abscisses porté par AB, écrire $A(-a;0)$, $B(a;0)$ avec $a>0$. Alors
\[
MA^2-MB^2=(x+a)^2+y^2-(x-a)^2-y^2=4ax.
\]
L'égalité des distances équivaut à $x=0$, droite passant par I et perpendiculaire à AB. Cela démontre les deux sens de la caractérisation.
\section{Distance minimale à une droite}
Si H est le projeté orthogonal de M sur une droite d, pour tout $P\in d$,
\[
MP^2=MH^2+HP^2\ge MH^2.
\]
Donc $MP\ge MH$, avec égalité uniquement si $P=H$. La perpendiculaire réalise la plus courte distance de M à la droite. Dans le cas de l'axe des abscisses, $H(x_M;0)$ et cette distance vaut $|y_M|$.
\section{Une preuve par choix de coordonnées}
Soit un rectangle ABCD non dégénéré. Choisir $A(0;0)$, $B(a;0)$, $C(a;b)$, $D(0;b)$, avec $a,b>0$. Les diagonales ont toutes deux pour carré de longueur $a^2+b^2$ ; elles ont donc la même longueur. Leurs milieux sont tous deux $(a/2;b/2)$ : elles se coupent en leur milieu. Le calcul est général car a et b représentent toutes les dimensions positives.
\begin{attention}
Un dessin fournit une représentation. Les mesures à la règle donnent des observations approchées ; elles ne prouvent pas une égalité exacte. En calculant, vérifier que le repère satisfait bien les hypothèses des formules utilisées.
\end{attention}
\begin{retenir}
Distance : différence des coordonnées puis somme des carrés, dans un repère orthonormé. Milieu : moyenne des coordonnées. Un lieu géométrique se traduit par une condition nécessaire et suffisante sur $(x;y)$.
\end{retenir}
\end{document}
